% =============================================================================
%  Section III — Messaging Subsystem
% =============================================================================

\section{Messaging Subsystem}
\label{sec:messaging}

The messaging layer of \ac{DFaaS} must serve two distinct communication
patterns across a dynamically changing peer-to-peer overlay: (i)~\emph{broadcast
dissemination} of shared load state to all reachable agents, and
(ii)~\emph{directed exchanges} between specific pairs of agents for
time-sensitive control operations such as offload negotiation and backpressure
signaling.
Rather than conflating these patterns in a single channel, \ac{DFaaS} adopts a
\emph{two-plane architecture} layered on the \emph{libp2p} networking stack,
separating concerns at the transport level while exposing a unified,
transport-agnostic message vocabulary to the strategy layer above.

% -----------------------------------------------------------------------
\subsection{Common Message Vocabulary}
\label{subsec:msg-vocabulary}

All inter-agent messages in \ac{DFaaS} share a common vocabulary defined in the
\texttt{msgtypes} package.
Every message embeds a \texttt{MsgHeader} structure that carries three fields:
a type discriminator (\texttt{msg\_type}), the libp2p peer ID of the
originating node (\texttt{sender\_id}), and a wall-clock timestamp.
The type discriminator enables \emph{cheap type-peeking}: before full \ac{JSON}
decoding, the receiving agent deserializes only a minimal
\texttt{MsgEnvelope} wrapper — a two-field struct containing only the header
— to determine the message kind, avoiding the allocation and reflection costs of
full message deserialization for unrecognized types.

The vocabulary is divided into two functional groups.
\emph{Broadcast messages} travel over the shared \ac{GossipSub} topic~\cite{vyzovitis2020gossipsub}:
\texttt{MsgHeartbeat} announces a node's presence, \ac{HAProxy} endpoint (host
and port), and the list of its locally deployed function names at regular
intervals;
\texttt{MsgOverloadAlert} signals that one or more functions have exceeded a
resource threshold, reporting the affected function names alongside current \ac{CPU}
and \ac{RAM} utilization normalized to $[0, 1]$;
and \texttt{MsgFunctionEvent} propagates immediate function lifecycle changes —
deployment and removal — so that peers can update their routing tables without
waiting for the next heartbeat cycle.
\emph{Directed messages} are exchanged point-to-point over libp2p streams:
\texttt{MsgBackpressure} is a fire-and-forget instruction sent to a peer
forwarding excessive traffic for specific functions;
\texttt{MsgOffloadRequest} opens a request/response negotiation in which the
sender proposes an additional forwarding rate in requests per second for a named
function;
and \texttt{MsgOffloadResponse} conveys the capacity the remote peer can absorb,
including a \texttt{correlation\_id} field echoed verbatim from the initiating
request to enable correct matching when multiple concurrent negotiations are in
flight.

% -----------------------------------------------------------------------
\subsection{Broadcast Plane: GossipSub}
\label{subsec:broadcast}

The broadcast plane is built on \ac{GossipSub}~\cite{vyzovitis2020gossipsub},
an application-layer pub/sub protocol designed for unstructured, decentralized
\ac{P2P} overlays.
Pub/sub systems have long been recognized as a natural fit for loosely coupled
distributed architectures~\cite{eugster2003pubsub}, and \ac{GossipSub} in
particular provides the epidemic dissemination properties~\cite{kermarrec2007epidemic}
needed to propagate load-state messages efficiently across a mesh whose
membership changes as agents join and leave the network.
\ac{GossipSub} constructs a bidirectional mesh among subscribed peers on a
shared topic and supplements eager push propagation with lazy \texttt{IHAVE}/\texttt{IWANT}
control messages to the fanout set, achieving reliable delivery while remaining
resilient to peer failures through a score-based peer-selection mechanism that
demotes poorly-performing mesh members.
In \ac{DFaaS}, all agents subscribe to a single configurable topic; the
\ac{GossipSub} router handles forwarding transparently, so strategies can
broadcast a message with a single \texttt{MarshAndPublish} call and expect
delivery to all live peers within the propagation latency of the mesh.

A dedicated goroutine runs \texttt{RunReceiver}, which reads messages from the
subscription using a blocking \texttt{Next} call and dispatches each to a
\emph{pre-filter callback chain}.
The outermost filter, \texttt{MakeCommonCallback}, performs two tasks on every
incoming message regardless of the active strategy: it checks whether the message
type belongs to the common broadcast vocabulary (via a set membership test on
the \texttt{CommonBroadcastTypes} map), and if so, dispatches the decoded message
to the appropriate update method on the \texttt{CommonNodeTable} before forwarding
to the inner strategy callback.
This layered design ensures that shared peer state is always kept current
irrespective of which strategy is active, without requiring individual strategy
implementations to duplicate parsing or table-management logic.

% -----------------------------------------------------------------------
\subsection{Common Node Table}
\label{subsec:nodestbl}

The \texttt{CommonNodeTable} is a thread-safe, \ac{TTL}-expiring in-memory store
keyed by peer ID.
Each entry records the peer's \ac{HAProxy} endpoint, the list of function names
it can execute, and the most recent \ac{CPU} and \ac{RAM} utilization reported by
\texttt{MsgOverloadAlert}.
Entries are seeded by heartbeats and refined incrementally by overload alerts and
function events; the \ac{TTL} is set to three times the heartbeat interval,
allowing entries to survive two consecutive missed heartbeats before expiring.
Read operations return a snapshot with deep-copied function slices to prevent
data races between concurrent strategy consumers and the callback goroutine
writing updates.

Self-messages — a consequence of \ac{GossipSub}'s forwarding model in which
the publishing agent may receive its own messages — are filtered at the strategy
level rather than at the transport.
The \texttt{MsgHeader.SenderID} field carried in every message enables strategy
implementations to compare against the local peer's ID and discard messages they
originated, avoiding spurious self-updates to the node table.

% -----------------------------------------------------------------------
\subsection{Directed Plane: libp2p Streams}
\label{subsec:directed}

The directed messaging plane is exposed through the \texttt{DirectMessenger}
interface, which wraps the libp2p stream \ac{API} under the protocol identifier
\texttt{/dfaas/msg/1.0.0}.
The \texttt{Send} method implements a fire-and-forget pattern: a new stream is
opened to the target peer, the serialized payload is written, and the stream is
closed without awaiting a reply.
The \texttt{Request} method extends this with a half-close handshake — after
writing the request, the sender calls \texttt{CloseWrite} to signal that it has
finished writing, then reads the response from the same open stream.
Both operations are subject to a configurable per-call timeout injected through
\texttt{context.WithTimeout}.

Wire framing uses a 4-byte big-endian length prefix followed by the
\ac{JSON}-encoded payload, assembled into a single buffer and issued as one
\texttt{Write} call to avoid partial-write races on multiplexed streams.
Typed handlers are registered via \texttt{SetHandler(msgType,~HandlerFunc)}
before the agent begins accepting connections; upon accepting an incoming stream,
\texttt{handleStream} peeks at the \texttt{MsgEnvelope} to identify the message
type, dispatches to the registered handler, and — if the handler returns a
non-nil response — writes the response back on the same stream.
Handlers returning \texttt{nil} implement fire-and-forget semantics from the
receiver's perspective, which is the correct behavior for \texttt{MsgBackpressure}.
